\documentclass[10pt,a4paper]{amsart}
\usepackage[margin=19mm]{geometry}
\usepackage{amsmath,amssymb,mathtools}
\usepackage[scr=rsfs,cal=euler]{mathalfa}
\usepackage{xfrac}
\usepackage[unicode,hidelinks,bookmarks=false]{hyperref}
\newcommand{\ie}{i.e.\ }
\newcommand{\cf}{cf.\ }
\newcommand{\resp}{respectively\ }
\newcommand{\Subsection}{\subsection}
\providecommand{\nobreakdash}{\nobreak\hbox{-}}
\setlength{\emergencystretch}{2em}
\multlinegap0pt
\usepackage[shortlabels]{enumitem}
\usepackage{amssymb}
\usepackage{tikz}
\usepackage{xcolor}
\usepackage{float}
\usepackage[normalem]{ulem}
\newtheorem{theorem}{Theorem}
\title{Collinear Interior Lattice Points in Triangles Satisfying $B(T)\in\{4,5\}$}
\author{Jonathan Sakunkoo, Annabella Sakunkoo,, Dana Paquin}
\date{Source: arXiv:2608.00889v1 (CC BY 4.0)}
\begin{document}
\maketitle

\noindent\emph{Source and licence.} Collinear Interior Lattice Points in Triangles Satisfying $B(T)\in\{4,5\}$. Version arXiv:2608.00889v1 (CC BY 4.0), distributed by arXiv under the Creative Commons Attribution 4.0 International licence, http://creativecommons.org/licenses/by/4.0/, as recorded in the arXiv licence metadata for this exact version and shown on the abstract page https://arxiv.org/abs/2608.00889v1. Full licence notice: \texttt{licenses/arxiv-2608.00889-CC-BY-4.0.txt}.\par\smallskip

\noindent\textit{Editorial scope.} Counted unit: the article's theorem five:label (source0.tex lines 227--230). The article's complete original proof of it is delivered with it (source0.tex lines 232--250, one \texttt{proof} environment of 19 lines). No mathematics is written, repaired or re-derived: the statement and the proof are the article's own lines, in the article's own notation, and no retained line is substituted or re-flowed.  No further source-local context is needed: every cross-reference of every retained line resolves inside the two delivered spans.  Nothing was omitted from any retained span, so the package carries no omission record.  No retained line contains a citation, so the wrapper carries no bibliography and no bibliography entry.  No retained line contains a figure environment, an image-inclusion command or drawing code, so no image is delivered and the package declares no asset dependency.  Editorial formatting only, in addition: an AMS \texttt{amsart} preamble that restates the article's own declarations and macros used by the retained lines, namely the environments \texttt{theorem} on separate counters, together with the standard packages those lines need.  The retained environments print in this document as Theorem 1 in order of appearance, whereas the article prints the same items with the numbering of its own full document.  The complete native source of the article as distributed in the arXiv e-print for this exact version is bundled unchanged as \texttt{sources/206575/source0.tex}.
\begin{theorem}
\label{five:label}
The number $5$ is $B4$-collinear. 
\end{theorem}

\begin{proof}
When $k=5$, the height $b$ must be $12$. 

By definition, $a$ is valid iff either $\gcd(a,12)=2$ and $\gcd(a-1,12)=1$, or $\gcd(a,12)=1$ and $\gcd(a-1,12)=2$. 

\textbf{Case 1.} In the former case, $a$ is even and $\gcd(a/2,6)=\gcd(a-1,3)=1$, so an even $a$ is valid iff $\gcd(a/2,6)=\gcd(a-1,3)=1$. Reducing modulo $12$, an even $a$ is in $\{0,2,4,6,8,10\}$; those satisfying $\gcd(a/2,6)=1$ are $2$ and $10$. Additionally, requiring $\gcd(a-1,3)=1$ shows that $a$ belongs to this case if and only if $a\equiv2\pmod{12}$. 

\textbf{Case 2.} In the latter case, $a$ is odd and $\gcd(a,3)=\gcd((a-1)/2,6)=1$. Thus, an odd $a$ is valid iff $\gcd(a,3)=\gcd((a-1)/2,6)=1$. Reducing modulo $12$, an odd $a$ is in $\{1,3,5,7,9,11\}$; those satisfying $\gcd(a,3)=1$ are $\{1,5,7,11\}$, and the only one $\gcd((a-1)/2,6)=1$ is $11$. 

Now, we have showed that the first case holds if and only if $a\equiv2\pmod{12}$, and the second case holds if and only if $a\equiv11\pmod{12}$. Recalling that $a$ is valid iff either the first or second case holds, it follows that $a$ is valid iff \[(a\bmod12)\in\{2,11\}.\] 

Note that the single non-vertex edge lattice point is at the midpoint of one of the edges; when $a\equiv2\pmod{12}$, it is at $(a/2,6)$, while when $a\equiv11\pmod{12}$ it is at $((a+1)/2,6)$. 

\textbf{Case 1.} Consider $a\equiv2\pmod{12}$; say $a=2+12m$ where $m$ is an integer. The open (i.e., excludes endpoints) line segment between $(1,0)$ and $(a/2,6)=(1+6m,6)$ lies entirely within the triangle's interior as it is a median of the triangle and contains exactly $5$ points as \[\gcd(1+6m-1,6-0)-1=6-1=5;\] these are all the interior points, and they are all collinear. 

\textbf{Case 2.} Consider $a\equiv11\pmod{12}$; say $a=11+12m$ where $m$ is an integer. The open median of the triangle from $(0,0)$ to $((a+1)/2,6)=(6+6m,6)$ lies entirely in the triangle and contains $5$ points as \[\gcd(6+6m,6)-1=6-1=5,\] which is all of the interior points - they are all collinear. 

Therefore, if $k=5$, all valid triangles have all of their interior points collinear. Then $5$ is $B4$-collinear. 
\end{proof}

\end{document}
